% !TeX program = xelatex

\documentclass[11pt]{exam}

%---enable russian----

\usepackage[utf8]{inputenc}
%\usepackage[russian]{babel}
\usepackage[english,main=russian]{babel}


\usepackage{cmap}
\usepackage{fontspec}
\usepackage{ulem}

\setmainfont{Roboto}
\setsansfont{Roboto}
\setmonofont{Roboto Mono}


\usepackage[margin=0.73in]{geometry}
%\usepackage[top=1in, bottom=1in, left=1in, right=1in]{geometry}

\usepackage{graphicx}
\usepackage{url}
\usepackage{latexsym}
\usepackage{amscd,amsmath,amsthm}
\usepackage{mathtools}
\usepackage{amsfonts}
\usepackage{amssymb}
\usepackage[dvipsnames]{xcolor}
\usepackage{hyperref}

\usepackage{algorithmicx, enumitem, algpseudocode, algorithm, caption}
\usepackage{tikz}
\usetikzlibrary{arrows,decorations.pathmorphing,backgrounds,calc}
\usetikzlibrary{chains, matrix, positioning, scopes, patterns, shapes}
\usepackage{pgfplots, subfigure}
\usetikzlibrary{automata}

\usepackage{fontawesome5}


\usepackage[most]{tcolorbox}

\usepackage[backend=biber,firstinits=true,hyperref=true,style=numeric-comp]{biblatex}

\definecolor{text-color}{RGB}{0,0,0} %%% text color
\definecolor{subsection-color}{RGB}{40, 166, 103} %%% subsection color
\definecolor{title-and-section-color}{RGB}{61, 147, 221} %%% title and section color
\definecolor{bg-color}{RGB}{254, 255, 244} %%% background color
\definecolor{struct-color}{RGB}{224, 62, 73} %%% structure

\definecolor{color-cont-bg}{RGB}{204, 239, 171}
\definecolor{block-green-color-bg}{RGB}{204, 239, 171}

\definecolor{plot-red-color}{RGB}{224, 62, 73}
\definecolor{plot-green-color}{RGB}{77, 158, 99}
\definecolor{plot-blue-color}{RGB}{61, 147, 221}

\definecolor{box-red-color}{RGB}{224, 62, 73}
\definecolor{box-green-color}{RGB}{77, 158, 99}
\definecolor{box-blue-color}{RGB}{61, 147, 221}

\bibliography{../biblio}

%%%%%%%%%%%%%%%%%%%%%
% Handling comments and versions %%%
%%%%%%%%%%%%%%%%%%%%%

%\renewcommand{\comment}[1]{\texttt{[#1]}}


%%%%%%%%%%%%%%%%%%%%%%%%%%%
%% THEOREMS
%%%%%%%%%%%%%%%%%%%%%%%%%%%

\newtheorem{theorem}{Theorem}[section]
\newtheorem{axiom}[theorem]{Axiom}
\newtheorem{conclusion}[theorem]{Conclusion}
\newtheorem{condition}[theorem]{Condition}
\newtheorem{conjecture}[theorem]{Conjecture}
\newtheorem{corollary}[theorem]{Corollary}
\newtheorem{criterion}[theorem]{Criterion}
\newtheorem{definition}[theorem]{Definition}
\newtheorem{lemma}[theorem]{Lemma}
\newtheorem{notation}[theorem]{Notation}
\newtheorem{proposition}[theorem]{Proposition}


\theoremstyle{definition}
\newtheorem{problem}{Problem}


\newcommand{\nc}{\newcommand}
\nc{\eps}{\varepsilon}
\nc{\RR}{{{\mathbb R}}}
\nc{\CC}{{{\mathbb C}}}
\nc{\FF}{{{\mathbb F}}}
\nc{\NN}{{{\mathbb N}}}
\nc{\ZZ}{{{\mathbb Z}}}
\nc{\PP}{{{\mathbb P}}}
\nc{\QQ}{{{\mathbb Q}}}
\nc{\UU}{{{\mathbb U}}}
\nc{\OO}{{{\mathbb O}}}
\nc{\EE}{{{\mathbb E}}}

\newcommand{\val}{\operatorname{val}}
\newcommand{\wt}{\ensuremath{\mathit{wt}}}
\newcommand{\Id}{\ensuremath{I}}
\newcommand{\transpose}{\mkern0.7mu^{\mathsf{ t}}}
\newcommand*{\ScProd}[2]{\ensuremath{\langle#1\mathbin{,}#2\rangle}} %Scalar Product
\renewcommand{\char}{\ensuremath{\mathsf{char}}}

\DeclareMathOperator{\Vol}{Vol}

\newcommand*{\structure}[1]{\textcolor{struct-color}{#1}}

%\pretolerance=1000

%%%%%%%%%%%%%%%%%%%%%%%%%%%%%%%%
%%%%%%%%%%%%%%%%%%%%%%%%%%%%%%%%
%% DOCUMENT STARTS
%%%%%%%%%%%%%%%%%%%%%%%%%%%%%%%%
%%%%%%%%%%%%%%%%%%%%%%%%%%%%%%%%


\begin{document}	
	{\noindent
		\textsc{БФУ им. И. Канта -- Компьютерный практикум по криптографии на эллиптических кривых }\\[5pt]
		Преподаватель {С. Новоселов}   \hfill{Осень 2025\\}
	\hrule
	\begin{center}
		{\LARGE\textbf{
				Лекция 2. Групповой закон \\[5pt]
		}}
		
	\end{center}
	\hrule \vspace{5mm}
	
	\thispagestyle{empty}
	
	\section{Определение}
	
	Пусть~$E/K : y^2 = x^3 + Ax + B$ -- эллиптическая кривая над полем~$K$. На множестве точек эллиптической кривой (плюс бесконечно удалённая точка) можно задать групповую операцию. Для кривых над полем действительных чисел~$\mathbb{R}$ это можно сделать следующим образом (см. рисунок~\ref{fig:group:law}). Пусть~$P_1 = (x_1, y_1)$ и~$P_2 = (x_2, y_2)$ -- точки эллиптической кривой.
	\begin{figure}[h!]
		\caption{Групповой закон над~$\mathbb{R}$}
		\label{fig:group:law}
		\centering
		\begin{tikzpicture}[scale=0.85]
			\begin{axis}[
				xmin=-4,
				xmax=5,
				xtick=\empty,
				ytick=\empty,
				ymin=-5,
				ymax=5,
				xlabel={$x$},
				ylabel={$y$},
				scale only axis,
				axis lines=middle,
				style={thick},
				domain=-2.279018:3,      
				samples=201,
				smooth,   
				clip=false,
				axis equal image=true,
				]
				\addplot[color=plot-blue-color] {sqrt(x^3-3*x+5)} node[right] {$E$};
				\addplot[color=plot-blue-color] {-sqrt(x^3-3*x+5)};
				\addplot[color=plot-red-color] coordinates {(-3, -0.0890722)(3, 3.06557)} node[right] {$\ell$};
				\addplot[color=plot-green-color] coordinates {(1.9, 3.5)(1.9, -3.5)};
				\draw [fill=black] (axis cs: 0.65, 1.83) circle (2pt);
				\draw[color=black] (axis cs: 1.2, 1.3) node [left]{$P_2$};
				\draw [fill=black] (axis cs: -2.26, 0.3) circle (2pt);
				\draw[color=black] (axis cs: -2.3, 0.3) node [left]{$P_1$};
				\draw [fill=black] (axis cs: 1.9, 2.5) circle (2pt);
				\draw[color=black] (axis cs: 1.9, 2.8) node [left]{$P_3'$};
				\draw [fill=black] (axis cs: 1.9, -2.5) circle (2pt);
				\draw[color=black] (axis cs: 1.9, -2.8) node [left]{$P_3$};
			\end{axis}
		\end{tikzpicture}
	\end{figure}
	Проведём прямую через данные точки, возьмём точку пересечения прямой с эллиптической кривой и отразим её относительно оси $x$, получив точку~$P_3 = \left(x_3, y_3\right)$. Решая систему уравнений, состоящую из уравнения кривой и уравнения прямой через две точки, получаем
	\begin{equation*}
		\begin{split}
			x_3 &= m^2 - x_1 - x_2, \\
			y_3 &= m\left( x_1 - x_3 \right) - y_1, \\
			m &= \frac{y_2 - y_1}{x_2 - x_1}.
		\end{split}
	\end{equation*}
	Может получиться так, что точки пересечения нет, как на Рис.~\ref{fig:group:law:inf}. В этом случае полагаем~$P_3 = \mathcal{O}$.
	\begin{figure}[h!]
		\centering
		\caption{Групповой закон над $\mathbb{R}$. Случай~$x_1 = x_2$.}
		\label{fig:group:law:inf}
		\begin{tikzpicture}[scale=0.85]
			\begin{axis}[
				xmin=-4,
				xmax=5,
				xtick=\empty,
				ytick=\empty,
				ymin=-5,
				ymax=5,
				xlabel={$x$},
				ylabel={$y$},
				scale only axis,
				axis lines=middle,
				style={thick},
				domain=-2.279018:3,      
				samples=201,
				smooth,   
				clip=false,
				axis equal image=true,
				]
				\addplot[color=plot-blue-color] {sqrt(x^3-3*x+5)} node[right] {$E$};
				\addplot[color=plot-blue-color] {-sqrt(x^3-3*x+5)};
				\addplot[color=plot-red-color] coordinates {(-1.0, 3.65)(-1.0, -3.65)};
				\draw [fill=black] (axis cs: -1.0, 2.65) circle (2pt);
				\draw[color=black] (axis cs: -1.0, 3.0) node [left]{$P_1$};
				\draw [fill=black] (axis cs: -1.0, -2.65) circle (2pt);
				\draw[color=black] (axis cs: -1.0, -3.0) node [left]{$P_2$};
			\end{axis}
		\end{tikzpicture}
	\end{figure}

	В случае, если $P_1 = P_2$, то вместо прямой, проходящей через две точки, берётся касательная к кривой в точке~$P_1$. В этом случае координаты точки $P_3$ получаются следующими:
	\begin{equation*}
		\begin{split}
			x_3 &= m^2 - 2 x_1, \\
			y_3 &= m\left( x_1 - x_3 \right) - y_1,\\
			m   &= \frac{3x_1^2 + A}{2 y_1}. \\ 
		\end{split}
	\end{equation*}
	\begin{figure}[h!]
			\caption{Групповой закон над $\mathbb{R}$. Удвоение точек.}
			\label{fig:group:law:dbl}
		\centering
		\begin{tikzpicture}[scale=0.85]
			\begin{axis}[
				xmin=-4,
				xmax=5,
				xtick=\empty,
				ytick=\empty,
				ymin=-5,
				ymax=5,
				xlabel={$x$},
				ylabel={$y$},
				scale only axis,
				axis lines=middle,
				style={thick},
				domain=-2.279018:3,      
				samples=201,
				smooth,   
				clip=false,
				axis equal image=true,
				]
				\addplot[color=plot-blue-color] {sqrt(x^3-3*x+5)} node[right] {$E$};
				\addplot[color=plot-blue-color] {-sqrt(x^3-3*x+5)};
				\addplot[color=plot-red-color] {2.65 + ((-3.0) + 3*(-1.2)^2)/(2*2.65) * (x - (-1.2))}  node[right] {$\ell$};
				\addplot[color=plot-green-color] coordinates {(2.5, 4.5)(2.5, -4.5)};
				\draw [fill=black] (axis cs: -1.2, 2.65) circle (2pt);
				\draw[color=black] (axis cs: -1.0, 3.0) node [left]{$P_1$};
				\draw [fill=black] (axis cs: 2.5, 3.6) circle (2pt);
				\draw[color=black] (axis cs: 2.5, 3.9) node [left]{$P_3'$};
				\draw [fill=black] (axis cs: 2.5, -3.6) circle (2pt);
				\draw[color=black] (axis cs: 2.5, -3.9) node [left]{$P_3$};
			\end{axis}
		\end{tikzpicture}
	\end{figure}
	
	Можно показать, что заданный таким способом групповой закон удовлетворяет всем аксиомам аддитивной абелевой группы.
	\begin{tcolorbox}[colframe=title-and-section-color!120, colback=title-and-section-color!5, title=Теорема, center title]
		%Операция сложения на $E$ обладает свойствами:
		\begin{enumerate}
			\item $P_1 + P_2 = P_2 + P_1$ \hfill \textit{(коммутативность)}
			
			\item $P + \mathcal{O} = P$ $\forall P \in E$ \hfill \textit{($\exists$ нейтральный элемент)}
			
			\item $\forall P \in E$ $\exists P' \in E: P + P' = \mathcal{O}$ \hfill \textit{($\exists$ обратный элемент)}
			
			\item $\left( P_1 + P_2 \right) + P_3 = P_1 + \left( P_2 + P_3 \right)$ \hfill \textit{(ассоциативность)}
		\end{enumerate}
	\end{tcolorbox}
	Более того, поскольку в формулах используются только операции сложения, умножения и деления, то они работают и для любого поля. А не только для~$\mathbb{R}$. Для криптографических приложений наиболее интересен случай~$K = \mathbb{F}_q$.
	
	\section{Быстрое скалярное умножение}
	Сложение точки самой собой несколько раз называется скалярным умножением, обозначается следующим образом.
	\[
	P \to \left[ k \right] \cdot P = \underbrace {P + P +  \ldots  + P}_{k{\text{-раз}}}
	\]
	
	Если складывать точку с самой собой наивно, то потребуется~$k-1$ сложений. Но для целей криптографии такой способ не подходит, так как в криптографических схемах используются большие значения~$k$. Есть более эффективный метод, представленный в Алгоритме~\ref{alg:scalar_mult}.
	
\begin{algorithm}
	\caption{Бинарный метод быстрого умножения}\label{alg:scalar_mult}
	\begin{algorithmic}
		\Require $k \in \mathbb{Z}, P \in E(\mathbb{F}_q)$
		\Ensure $[k] P$
		\State $k = \sum\limits_{j = 0}^{\ell - 1} {{k_j}{2^j}},\quad {k_j} \in \left\{ {0,1} \right\}$\
		\State $Q \gets \mathcal{O}$
		\For{$j = \ell - 1$ \textbf{to} $0$ \textbf{by} $-1$}
			\State $Q \gets \left[ 2 \right]Q$
			\If{${k_j} = 1$}
				\State $Q \gets Q + P$
			\EndIf
		\EndFor~\\
		\Return $Q$
	\end{algorithmic}
\end{algorithm}
Он уже требует только~$\mathcal{O}(\log k)$ сложений точек.

\section{Протокол Диффи-Хэллмана}
Группового закона и быстрого алгоритма скалярного умножения уже достаточно для определения криптографических схем. На следующем схеме представлен протокол обмена ключами Диффи-Хеллмана, позволяющий двум абонентам выработать общий секретный ключ.

\begin{center}

\begin{tikzpicture}[x=0.75pt,y=0.75pt,yscale=-1,xscale=1]
	%Straight Lines [id:da0832972187828418] 
	\draw    (80,35) -- (185,35) ;
	\draw [shift={(188,35)}, rotate = 540] [fill={rgb, 255:red, 0; green, 0; blue, 0 }  ][line width=0.08]  [draw opacity=0] (5.36,-2.57) -- (0,0) -- (5.36,2.57) -- cycle    ;
	\draw [shift={(77,35)}, rotate = 360] [fill={rgb, 255:red, 0; green, 0; blue, 0 }  ][line width=0.08]  [draw opacity=0] (5.36,-2.57) -- (0,0) -- (5.36,2.57) -- cycle    ;
	%Straight Lines [id:da17602726868869878] 
	\draw    (80,70) -- (185,70) ;
	\draw [shift={(188,70)}, rotate = 540] [fill={rgb, 255:red, 0; green, 0; blue, 0 }  ][line width=0.08]  [draw opacity=0] (5.36,-2.57) -- (0,0) -- (5.36,2.57) -- cycle    ;
	%Straight Lines [id:da5151276494097119] 
	\draw    (80,105) -- (185,105) ;
	\draw [shift={(79,105)}, rotate = 360] [fill={rgb, 255:red, 0; green, 0; blue, 0 }  ][line width=0.08]  [draw opacity=0] (5.36,-2.57) -- (0,0) -- (5.36,2.57) -- cycle    ;
	
	% Text Node
	\draw (50,24) node [anchor=north west][inner sep=0.75pt]   [align=left] { \structure{{\Large\faUserSecret}} };
	\draw (200,24) node [anchor=north west][inner sep=0.75pt]   [align=right] {\structure{{\Large\faCat}}};
	\draw (102,15) node [anchor=north west][inner sep=0.75pt]   [align=center] {$P \in E(\mathbb{F}_q)$};
	% Text Node
	\draw (35,62) node [anchor=north west][inner sep=0.75pt]   [align=left] { \structure{$a \in \mathbb{N}$} };
	\draw (120,50) node [anchor=north west][inner sep=0.75pt]   [align=left] {$[a] P$};
	% Text Node
	\draw (120,86) node [anchor=north west][inner sep=0.75pt]   [align=left] {$[b]P$};
	\draw (190,98) node [anchor=north west][inner sep=0.75pt]   [align=right] { \structure{$b \in \mathbb{N}$} };
	% Text Node
	\draw (-40,130) node [anchor=north west][inner sep=0.75pt]   [align=left] {\structure{\boxed{[ab]P = [a]([b]P)}}};
	\draw (190,130) node [anchor=north west][inner sep=0.75pt]   [align=left] {\structure{\boxed{[ab]P = [b]([a]P)}}};
\end{tikzpicture}
\end{center}
В схеме каждый абонент выбирает свой секретный ключ -- число $a$ (или~$b$), а затем пересылает по открытому каналу значение~$[a]P$ (или~$[b]P$), которое представляют собой открытый ключ. Затем оба абонента по открытым ключам и своему секретному ключу вычисляют общий секретный ключ. Безопасность схемы строится на том факте, что задача вычисления числа $a$ зная только генератор группы~$P$ и точку~$[a]P$, является сложной задачей. Данная задача называется \structure{задачей вычисления дискретного логарифма (DLOG)}. На сегодняшний день не известно классического (не квантового) алгоритма, который в общем случае решал бы данную задачу за полиномиальное время.

Заметим, что на практике схемы шифрования с открытым ключом не используются для шифрования большого количества информации из-за медленной (по сравнению с симметричными шифрами) скорости работы. Вместо этого используется гибридный подход -- сначала вырабатывается общий секретный ключ с использованием протокола обмена ключами. А затем данный ключ применяется для шифрования данных с помощью симметричного шифра.

\section{Атака ``человек посередине''}
Протокол Диффи-Хелмана уязвим к атаке ''человек посередине'', в которой имеется третья сторона, которая имеет возможность перехватывать и подменять открытые сообщения абонентов. В качестве такой третьей стороны может выступать провайдер интернета или владелец точки доступа. Схема атаки представлена на следующем рисунке.

\begin{center}
	\tikzset{every picture/.style={line width=0.75pt}} %set default line width to 0.75pt        

\begin{tikzpicture}[x=0.75pt,y=0.75pt,yscale=-0.98,xscale=0.98]
	%%%%%%%%%%%%%%%%%%%
	% Left side arrows
	%%%%%%%%%%%%%%%%%%%
	%Straight Lines [id:da17602726868869878] 
	\draw (145,130) -- (205,130);
	\draw [shift={(207,130)}, rotate = 180] [fill={rgb, 255:red, 0; green, 0; blue, 0 }  ][line width=0.08]  [draw opacity=0] (5.36,-2.57) -- (0,0) -- (5.36,2.57) -- cycle;
	
	%Straight Lines [id:da606368011362501]
	\draw    (145,165) -- (205,165) ;
	\draw [shift={(142,165)}, rotate = 360] [fill={rgb, 255:red, 0; green, 0; blue, 0 }  ][line width=0.08]  [draw opacity=0] (5.36,-2.57) -- (0,0) -- (5.36,2.57) -- cycle;
	
	%Straight Lines [id:da0586486120331462] 
	\draw    (145,195) -- (205,195) ;
	\draw [shift={(207,195)}, rotate = 180] [fill={rgb, 255:red, 0; green, 0; blue, 0 }  ][line width=0.08]  [draw opacity=0] (5.36,-2.57) -- (0,0) -- (5.36,2.57) -- cycle;
	
	%%%%%%%%%%%%%%%%%%%%
	% Right side arrows
	%%%%%%%%%%%%%%%%%%%%
	%Straight Lines [id:da5151276494097119] 
	\draw (320,130) -- (380,130) ;
	\draw [shift={(318,130)}, rotate = 360] [fill={rgb, 255:red, 0; green, 0; blue, 0 }  ][line width=0.08]  [draw opacity=0] (5.36,-2.57) -- (0,0) -- (5.36,2.57) -- cycle;
	
	%Straight Lines [id:da09168079881977365] 
	\draw    (320,165) -- (380,165) ;
	\draw [shift={(382,165)}, rotate = 180] [fill={rgb, 255:red, 0; green, 0; blue, 0 }  ][line width=0.08]  [draw opacity=0] (5.36,-2.57) -- (0,0) -- (5.36,2.57) -- cycle;
	
	%Straight Lines [id:da4782643167070846] 
	\draw    (320,195) -- (380,195) ;
	\draw [shift={(382,195)}, rotate = 180] [fill={rgb, 255:red, 0; green, 0; blue, 0 }  ][line width=0.08]  [draw opacity=0] (5.36,-2.57) -- (0,0) -- (5.36,2.57) -- cycle;
	
	%%%%%%%%%%%%%%%%%%%%%%%
	% Text nodes for Agent
	%%%%%%%%%%%%%%%%%%%%%%%
	% Text Node
	\draw (120,72) node [anchor=north west][inner sep=0.75pt]   [align=left] {\structure{{\Large\faUserSecret}}};
	
	% Text Node
	\draw (100,122) node [anchor=north west][inner sep=0.75pt]   [align=left] {\structure{$a \in \mathbb{N}$}};
	
	% Text Node
	\draw (52,156) node [anchor=north west][inner sep=0.75pt]   [align=left] {\structure{$k_a = [a] ([c] P)$}};
	
	% Text Node
	\draw (42,186) node [anchor=north west][inner sep=0.75pt]   [align=left] {\structure{$c_a=Enc_{k_a}(\text{\faCheese})$}};
	
	%%%%%%%%%%%%%%%%%%%%%%%%%%%%%%%%%
	% Text nodes for left side arrows
	%%%%%%%%%%%%%%%%%%%%%%%%%%%%%%%%%
	% Text Node
	\draw (162,110) node [anchor=north west][inner sep=0.75pt]   [align=left] {$[a] P$};
	
	% Text Node
	\draw (162,146) node [anchor=north west][inner sep=0.75pt]   [align=left] {$[c] P$};
	
	% Text Node
	\draw (170,180) node [anchor=north west][inner sep=0.75pt]   [align=left] {$c_a$};
	
	%%%%%%%%%%%%%%%%%%%%%%%%%%%%%%%%%%%
	% Text nodes for public parameters
	%%%%%%%%%%%%%%%%%%%%%%%%%%%%%%%%%%%
	% Text Node
	\draw (225,35) node [anchor=north west][inner sep=0.75pt]   [align=left] {\structure{\boxed{P \in E(\mathbb{F}_q)}}};
	
	%%%%%%%%%%%%%%%%%%%%%%%
	% Text nodes for Mouse
	%%%%%%%%%%%%%%%%%%%%%%%
	% Text Node
	\draw (245,115) node [anchor=north west][inner sep=0.75pt]   [align=center] {\includegraphics[width=2.5em]{mouse.png}};
	
	% Text Node
	\draw (245,155) node [anchor=north west][inner sep=0.75pt]   [align=center] {\structure{$c \in \mathbb{N}$}};
	
	% Text Node
	\draw (215,186) node [anchor=north west][inner sep=0.75pt]   [align=center] {\structure{$\text{\faCheese} = Dec_{k_a}(c_a)$}};
	
	% Text Node
	\draw (216,215) node [anchor=north west][inner sep=0.75pt]   [align=left] {\structure{$c_b = Enc_{k_b}(\text{\faCheese})$}};
	
	
	%%%%%%%%%%%%%%%%%%%%%%%%%%%%%%%%%%
	% Text Nodes for right side arrows
	%%%%%%%%%%%%%%%%%%%%%%%%%%%%%%%%%%
	% Text Node
	\draw (335,110) node [anchor=north west][inner sep=0.75pt]   [align=left] {$[b] P$};	
	
	% Text Node
	\draw (335,146) node [anchor=north west][inner sep=0.75pt]   [align=left] {$[c] P$};
	
	% Text Node
	\draw (340,180) node [anchor=north west][inner sep=0.75pt]   [align=left] {$c_b$};
	
	
	%%%%%%%%%%%%%%%%%%%%%%%%%%%%%%%%%%
	% Text Nodes for Cat
	%%%%%%%%%%%%%%%%%%%%%%%%%%%%%%%%%%
	% Text Node
	\draw (385,72) node [anchor=north west][inner sep=0.75pt]   [align=left] {\structure{{\Large\faCat}}};
	
	% Text Node
	\draw (385,122) node [anchor=north west][inner sep=0.75pt]   [align=left] {\structure{$b \in \mathbb{N}$}};
	
	% Text Node
	\draw (385,156) node [anchor=north west][inner sep=0.75pt]   [align=left] {\structure{$k_b = [b] ([c] P)$}};
	
	% Text Node
	\draw (385,186) node [anchor=north west][inner sep=0.75pt]   [align=left] {\structure{$\text{\faCheese} = Dec_{k_b}(c_b)$}};
	
\end{tikzpicture}

\end{center}
Третья сторона вырабатывает с каждым из абонентов свой секретный ключ и с помощью секретных ключей перешифровывает переписку, попутно читая её. 

%\section{Цифровая подпись}


\nocite{Blake1999}
\nocite{Washington2008}
\nocite{Menezes1993}
\nocite{HankersonMenezesVanstone2006}
\nocite{Silverman2009}
\printbibliography

\end{document}
